\documentclass[11pt]{article}

\usepackage[margin=1in]{geometry}
\usepackage{amsmath, amssymb}
\usepackage{booktabs}
\usepackage{graphicx}
\usepackage[hypertexnames=false]{hyperref}
\usepackage{enumitem}
\usepackage{float}
\setlist{leftmargin=*}
\setlength{\parskip}{0.35em}
\setlength{\parindent}{0pt}

\title{\textbf{Experimental Report}\\
Submodular Optimization and Greedy Approximation for Distributed Influence Maximization}
\author{Yixin Chen \and Bichi Zhang \and Yaoqin Ye}
\date{CS240: Algorithm Design and Analysis, Spring 2026}

\begin{document}

\maketitle

\begin{abstract}
This project studies influence maximization, a classical network optimization
problem that asks how to choose a small number of seed nodes so that a diffusion
process reaches as many nodes as possible. Following the CS240 project
requirements, we connect a modern data application to classical algorithms:
monotone submodular maximization, greedy approximation, lazy evaluation, and a
communication-constrained distributed variant. We implemented Independent
Cascade simulation, random/degree/PageRank baselines, centralized Monte Carlo
greedy, CELF lazy greedy, and a GreeDI-style local candidate sharing method.
Experiments were run on synthetic and small real-world benchmark graphs. The
results show that structural baselines are surprisingly strong on small graphs,
that CELF reduces the number of influence evaluations while preserving the same
algorithmic objective, and that distributed candidate sharing exposes a clear
tradeoff among solution quality, runtime, and communication budget.
\end{abstract}

\section{Background and Related Work}

Influence maximization is a standard problem in social networks, viral
marketing, public-health intervention, and misinformation analysis. Given a
network and a limited budget, the goal is to choose seed nodes that maximize
the expected number of eventually activated nodes. The problem is algorithmic
rather than purely empirical: the chosen seed set must solve a combinatorial
optimization problem under uncertainty.

Kempe, Kleinberg, and Tardos~\cite{kempe2003} introduced influence
maximization under the Independent Cascade (IC) and Linear Threshold (LT)
models. They showed that the problem is NP-hard, but that the expected spread
function is monotone and submodular. Nemhauser, Wolsey, and Fisher's classical
result~\cite{nemhauser1978} then implies that greedy maximization under a
cardinality constraint obtains a $(1-1/e)$ approximation for the exact
submodular objective.

The difficulty is that exact influence spread is expensive to compute. A
straightforward implementation estimates the spread of many candidate seed sets
using Monte Carlo simulation. Leskovec et al.~\cite{leskovec2007} proposed
CELF, a lazy greedy method that uses submodularity to avoid many repeated
marginal-gain evaluations. For large-scale or distributed settings,
Mirzasoleiman et al.~\cite{mirzasoleiman2013} proposed GreeDI, which runs
greedy locally on partitions and then merges local representatives. More recent
methods such as TIM/IMM~\cite{tang2014,tang2015} use reverse influence sampling
to scale influence maximization further; in this project, those methods are
treated as related work rather than full reproduction targets.

\section{Problem Formulation}

Let $G=(V,E,p)$ be a directed graph. Each edge $(u,v)$ has an activation
probability $p_{uv}\in[0,1]$. Given a seed budget $k$, the influence
maximization problem is
\[
    \max_{S\subseteq V,\ |S|\leq k} \sigma(S),
\]
where $\sigma(S)$ is the expected number of active nodes after diffusion starts
from seed set $S$.

Our implementation uses the Independent Cascade model. When a node $u$ first
becomes active, it gets one chance to activate each inactive out-neighbor $v$
with probability $p_{uv}$. Each activation attempt is independent, and the
process stops when no newly active nodes remain.

Under the IC model, $\sigma$ is monotone and submodular:
\[
    \sigma(A\cup\{v\})-\sigma(A)
    \geq
    \sigma(B\cup\{v\})-\sigma(B),
    \quad A\subseteq B,\ v\notin B.
\]
This diminishing-returns property is the reason greedy selection has a
theoretical guarantee and the reason CELF can safely treat old marginal gains
as upper bounds.

\section{Representative Method Analysis}

\subsection{Centralized Greedy}

The centralized greedy algorithm starts with $S_0=\emptyset$. At round $t$, it
chooses
\[
    v^\star =
    \arg\max_{v\in V\setminus S_t}
    \left(\widehat{\sigma}(S_t\cup\{v\})-\widehat{\sigma}(S_t)\right),
\]
where $\widehat{\sigma}$ is the Monte Carlo estimate of IC spread. The seed set
is updated as $S_{t+1}=S_t\cup\{v^\star\}$ until $|S|=k$.

If the graph has $n$ nodes and $m$ edges, one IC simulation costs
$O(n+m)$ in the worst case. With $R$ Monte Carlo simulations per estimate,
naive greedy costs approximately
\[
    O(k n R (n+m)),
\]
ignoring small savings from already selected nodes. This is feasible for small
course-project graphs but becomes expensive quickly.

\subsection{CELF Lazy Greedy}

CELF maintains a priority queue of cached marginal gains. Since submodularity
implies that a node's marginal gain cannot increase as the seed set grows, an
old marginal gain is an upper bound on its current gain. When the top node in
the queue was last evaluated for the current seed-set size, CELF accepts it.
Otherwise, it recomputes only that node's marginal gain and pushes it back.
The selected seed set targets the same greedy objective but typically requires
far fewer spread evaluations.

\subsection{GreeDI-Style Distributed Greedy}

For the distributed component, the node set is partitioned into $m$ local
parts. Worker $i$ runs greedy over its local candidate list and returns at most
$q$ candidates. The coordinator forms
\[
    C = \bigcup_{i=1}^{m} C_i
\]
and runs a final greedy selection over $C$ rather than all nodes. We test
$m\in\{2,4,8\}$ and $q\in\{k,2k,5k\}$. The communication proxy is the number of
transmitted candidates, approximately $mq$ but capped by partition sizes.

\section{Reproduction and Implementation}

The source code is organized as follows:
\begin{itemize}
    \item \texttt{src/simulation.py}: Independent Cascade simulation and Monte
    Carlo spread estimation.
    \item \texttt{src/baselines.py}: random, weighted out-degree, and PageRank
    seed selection.
    \item \texttt{src/greedy.py}: centralized greedy and CELF lazy greedy.
    \item \texttt{src/distributed.py}: random node partitioning and
    GreeDI-style candidate sharing.
    \item \texttt{src/graphs.py}: benchmark graph construction and
    weighted-cascade probability assignment.
    \item \texttt{src/experiments.py}: end-to-end experiment runner and CSV
    output.
\end{itemize}

Because exact reproduction of large influence maximization systems is beyond
the scope of the course project, we reproduced the core algorithmic claims at a
smaller scale: IC spread estimation, greedy submodular selection, CELF lazy
evaluation, and distributed local-candidate merging. The experiment runner used
six cases: Erdos--Renyi graphs with 40 and 70 nodes, Barabasi--Albert graphs
with 40 and 70 nodes, a Watts--Strogatz graph with 60 nodes, and Zachary's
Karate Club graph with 34 nodes. The seed budget was $k=5$ or $k=6$ depending
on the case.

Each algorithm's final seed set was re-evaluated using a common Monte Carlo
seed for the same graph case. This makes the reported spread more comparable
than the internal spread estimates used during selection.

\section{Experimental Evaluation}

\subsection{Overall Algorithm Comparison}

Table~\ref{tab:overall} summarizes the average results across all graph cases.
The reported ratio is relative to centralized greedy on the same graph.

\begin{table}[H]
\centering
\caption{Average performance across all benchmark cases.}
\label{tab:overall}
\begin{tabular}{lrrrr}
\toprule
Algorithm & Spread & Ratio & Runtime (s) & Evaluations \\
\midrule
Random & 16.58 & 0.93 & 0.00 & 1.0 \\
Degree & 24.42 & 1.37 & 0.00 & 1.0 \\
PageRank & 22.92 & 1.27 & 0.00 & 1.0 \\
Greedy & 18.08 & 1.00 & 0.06 & 283.5 \\
CELF & 18.92 & 1.04 & 0.02 & 112.0 \\
Best distributed average setting & 24.00 & 1.34 & 0.32 & 359.5 \\
\bottomrule
\end{tabular}
\end{table}

The most important observation is that the small-scale experiment does not show
a clean dominance of greedy over structural baselines. Degree and PageRank
perform very well on these graphs, especially because the weighted-cascade
probability rule favors high-degree structural positions. This is not a
contradiction of the greedy approximation theorem: the theorem compares greedy
to the optimal seed set for the exact submodular objective, while our
implementation uses noisy Monte Carlo estimates during selection and then
re-evaluates final seed sets with another random seed.

CELF behaves as expected: it obtains similar average spread to greedy but uses
fewer spread evaluations, reducing the average from 283.5 to 112.0 evaluations.
This directly demonstrates the benefit of lazy evaluation.

\begin{figure}[H]
\centering
\includegraphics[width=0.92\linewidth]{../outputs/influence_maximization/figures/estimated_spread.png}
\caption{Estimated IC spread by graph size and algorithm.}
\end{figure}

\subsection{Distributed Communication Tradeoff}

Table~\ref{tab:distributed} summarizes the distributed variants. The notation
$(m,q/k)$ means $m$ partitions and local candidate budget $q=(q/k)\cdot k$.

\begin{table}[H]
\centering
\caption{Average distributed performance under different partition and communication budgets.}
\label{tab:distributed}
\begin{tabular}{lrrrrr}
\toprule
Setting & Spread & Ratio & Runtime (s) & Evaluations & Candidates \\
\midrule
$(2,1)$ & 21.58 & 1.19 & 0.30 & 321.5 & 11.0 \\
$(2,2)$ & 21.00 & 1.16 & 0.64 & 591.5 & 22.0 \\
$(2,5)$ & 21.50 & 1.20 & 1.37 & 1021.5 & 49.0 \\
$(4,1)$ & 24.00 & 1.34 & 0.32 & 359.5 & 22.0 \\
$(4,2)$ & 21.50 & 1.18 & 0.61 & 599.8 & 43.0 \\
$(4,5)$ & 22.08 & 1.22 & 0.71 & 683.8 & 52.3 \\
$(8,1)$ & 22.00 & 1.21 & 0.42 & 430.5 & 43.0 \\
$(8,2)$ & 21.00 & 1.16 & 0.58 & 503.2 & 52.3 \\
$(8,5)$ & 21.92 & 1.22 & 0.51 & 503.2 & 52.3 \\
\bottomrule
\end{tabular}
\end{table}

The distributed results show the expected cost side of the tradeoff clearly:
larger local budgets generally increase the number of transmitted candidates
and spread evaluations. The quality side is noisier. In several cases,
distributed candidate sharing beats the centralized greedy baseline after fair
re-evaluation. This can happen because local greedy introduces diversity into
the candidate pool and because the centralized greedy baseline is selected
under noisy Monte Carlo estimates. The result should therefore be interpreted
as evidence that candidate sharing can be competitive on small graphs, not as a
claim that it theoretically dominates centralized greedy.

\begin{figure}[H]
\centering
\includegraphics[width=0.92\linewidth]{../outputs/influence_maximization/figures/quality_ratio_to_greedy.png}
\caption{Quality ratio relative to centralized greedy.}
\end{figure}

\begin{figure}[H]
\centering
\includegraphics[width=0.92\linewidth]{../outputs/influence_maximization/figures/transmitted_candidates.png}
\caption{Communication proxy: number of transmitted candidates.}
\end{figure}

\subsection{Runtime and Evaluation Counts}

Centralized greedy has a small wall-clock time on these graphs, but its
evaluation count grows with $k$ and the number of candidate nodes. CELF reduces
the number of evaluations by reusing marginal-gain upper bounds. Distributed
variants perform more total evaluations because each worker runs a local greedy
phase and the coordinator runs another greedy phase. This is acceptable only
when worker computations are parallelized or when full-graph candidate search
is impossible due to memory or ownership constraints.

\begin{figure}[H]
\centering
\includegraphics[width=0.92\linewidth]{../outputs/influence_maximization/figures/spread_evaluations.png}
\caption{Number of spread evaluations. CELF reduces evaluations compared with naive greedy.}
\end{figure}

\begin{figure}[H]
\centering
\includegraphics[width=0.92\linewidth]{../outputs/influence_maximization/figures/runtime_seconds.png}
\caption{Wall-clock runtime on the current benchmark set.}
\end{figure}

\section{Deviations, Limitations, and Optional Extensions}

The proposal mentioned larger synthetic graphs and possible real networks such
as Wiki-Vote or NetHEPT. The current reproduction uses smaller graphs so that
the whole experiment can run quickly in the course-project environment. This
choice keeps the implementation transparent, but it limits scalability claims.

Another limitation is Monte Carlo noise. The experiment uses a modest number of
simulations for speed, and distributed variants use fewer simulations during
their local selections. As a result, the selected seed set may be sensitive to
randomness. This explains why the fair re-evaluation sometimes ranks degree or
distributed variants above centralized greedy. A stronger final experiment
would increase the simulation count, average over multiple random seeds, and
report confidence intervals.

The implemented optional extension is the GreeDI-style communication-constrained
variant. Further possible extensions include deterministic partition strategies,
larger real datasets, RR-set based maximum coverage, and parallel worker
execution to measure true distributed speedup.

\section{Conclusion}

This project demonstrates how a modern network problem can be understood using
classical algorithmic ideas. Influence maximization under IC diffusion gives a
monotone submodular objective; greedy selection provides a theoretical
approximation guarantee; CELF turns submodularity into a practical acceleration;
and GreeDI-style candidate sharing connects the problem to distributed
optimization under communication limits. The empirical results are intentionally
small-scale, but they reproduce the central algorithmic mechanisms and expose
the main practical issue: influence maximization quality depends not only on
the approximation algorithm, but also on spread-estimation accuracy, graph
structure, and communication constraints.

\begin{thebibliography}{9}

\bibitem{kempe2003}
D. Kempe, J. Kleinberg, and E. Tardos.
\newblock Maximizing the spread of influence through a social network.
\newblock \emph{KDD}, 2003.

\bibitem{nemhauser1978}
G. L. Nemhauser, L. A. Wolsey, and M. L. Fisher.
\newblock An analysis of approximations for maximizing submodular set functions.
\newblock \emph{Mathematical Programming}, 1978.

\bibitem{leskovec2007}
J. Leskovec, A. Krause, C. Guestrin, C. Faloutsos, J. VanBriesen, and N. Glance.
\newblock Cost-effective outbreak detection in networks.
\newblock \emph{KDD}, 2007.

\bibitem{mirzasoleiman2013}
B. Mirzasoleiman, A. Karbasi, R. Sarkar, and A. Krause.
\newblock Distributed submodular maximization: Identifying representative elements in massive data.
\newblock \emph{NeurIPS}, 2013.

\bibitem{tang2014}
Y. Tang, X. Xiao, and Y. Shi.
\newblock Influence maximization: Near-optimal time complexity meets practical efficiency.
\newblock \emph{SIGMOD}, 2014.

\bibitem{tang2015}
Y. Tang, Y. Shi, and X. Xiao.
\newblock Influence maximization in near-linear time: A martingale approach.
\newblock \emph{SIGMOD}, 2015.

\end{thebibliography}

\end{document}
